\chapter{مقدمه}
\section{بیان مسئله}
بهینه‌سازی (Optimization) علمی است که به یافتن بهترین جواب ممکن از میان مجموعه‌ای از جواب‌های شدنی می‌پردازد. یک مسئله بهینه‌سازی را می‌توان به صورت زیر مدل کرد:
\begin{equation}
	\min_{x \in \mathcal{X}} f(x)
	\label{eq:opt-problem}
\end{equation}
که در آن \( f : \mathbb{R}^n \rightarrow \mathbb{R} \) تابع هدف و \( \mathcal{X} \subseteq \mathbb{R}^n \) مجموعه شدنی است. در بسیاری از کاربردهای عملی مانند یادگیری ماشین، کنترل بهینه و طراحی مهندسی، توابع هدف غیرخطی و محدودیت‌های پیچیده‌ای وجود دارد که حل تحلیلی را ناممکن می‌سازد. بنابراین نیاز به روش‌های عددی تکراری احساس می‌شود.

\section{اهمیت و کاربردها}
روش گرادیان کاهشی (Gradient Descent) به دلیل سادگی و کارایی در مسائل با ابعاد بالا، به پرکاربردترین الگوریتم بهینه‌سازی در یادگیری ماشین تبدیل شده است. از جمله کاربردهای آن می‌توان به تنظیم پارامترهای شبکه‌های عصبی، رگرسیون خطی و لجستیک، و بهینه‌سازی توابع محدب اشاره کرد. در این پروژه تمرکز اصلی بر روی نسخه‌ی نامقید گرادیان کاهشی با گام ثابت خواهد بود.

\section{ساختار پایان‌نامه}
این پایان‌نامه در چهار فصل سازماندهی شده است: 
\begin{itemize}
	\item در **فصل دوم** فرمول‌های پایه‌ای شامل فرمول طولانی در محیط align، فرمول چندضابطه‌ای و فرمول نرم برداری ارائه می‌شود.
	\item **فصل سوم** به معرفی الگوریتم گرادیان کاهشی اختصاص دارد؛ شبه‌کد الگوریتم، فلوچارت مربوطه و یک جدول مقایسه‌ای از روش‌های بهینه‌سازی در این فصل آورده می‌شود. همچنین یک شکل با دو زیرشکل برای نمایش رفتار الگوریتم ترسیم شده است.
	\item در **فصل چهارم** نتایج و دستاوردهای پروژه جمع‌بندی شده و پیشنهاداتی برای کارهای آینده ارائه می‌گردد. همچنین یک گراف تعاملی برای نشان دادن ارتباط فصول طراحی شده است.
\end{itemize}